% Curs 10 - Pagina 8 (6-XII-1994).
% Fragment pentru % Curs 10 - Pagina 8 (6-XII-1994).
% Fragment pentru % Curs 10 - Pagina 8 (6-XII-1994).
% Fragment pentru % Curs 10 - Pagina 8 (6-XII-1994).
% Fragment pentru \input{courses/pages/curs10-pagina8-ro}.
% Fara preambul si fara intrerupere fortata de pagina.
% Continuarea demonstratiei lemei 2 de pe pagina 7.
% Nota de fidelitate: manuscrisul pare sa aiba "a in A_0",
% desi A_0 este o submultime de clase aH. Lectura este pastrata.

\subsection*{Curs 10 - Pagina 8}

\noindent $a\in A_0\ \Longleftrightarrow\ haH=aH$, $(\forall)\,h\in H$
\[
\Longleftrightarrow\ a^{-1}haH=H,\quad(\forall)\,h\in H
\]
\[
\Longleftrightarrow\ a^{-1}Ha=H
\Longleftrightarrow\ H=aHa^{-1}
\Longleftrightarrow\ a\in N_G(H).
\]

\noindent Deci
\[
|A_0|=\left|\frac{N_G(H)}{H}\right|=[N_G(H):H],
\]
\noindent iar $|A|=[G:H]$.

\bigskip
\noindent\underline{Corolar}. Aceleaşi ipoteze ca la lema 2 şi
$p\mid[G:H]$.

\noindent Atunci $H\subsetneq N_G(H)$.

\bigskip
\noindent\fbox{\textbf{Dem. T. lui Sylow 1}}\quad $p^i\mid |G|$.

\medskip
\noindent Inducţie după $i$.

\noindent $i=1$: O.K. din T. lui Cauchy.

\medskip
\noindent Presupunem $1\leq i<t$ şi că există $H\leq G$ a.î.
$|H|=p^i$.

\noindent Demonstrăm că există $K\leq G$, $|K|=p^{i+1}$,
şi $H\trianglelefteq K$.

\medskip
\noindent $p\mid[G:H]$, deoarece
\[
|G|=p^t m=[G:H]\,|H|=[G:H]p^i.
\]
\noindent $i<t\ \Longrightarrow\ p^{t-i}\mid[G:H]
\ \Longrightarrow\ p\mid[G:H]$.

\medskip
\noindent Din corolar, $H\subsetneq N_G(H)$.
Considerăm $N_G(H)/H$.

\medskip
\noindent Afirm că
\[
p\mid\left|\frac{N_G(H)}{H}\right|=[N_G(H):H],
\]
\noindent deoarece $p\mid[G:H]$ şi
\[
[N_G(H):H]\equiv[G:H]\pmod p.
\]

\noindent Din T. lui Cauchy, $(\exists)\,\hat a\in N_G(H)/H$
a.î. $\operatorname{ord}(\hat a)=p$.

\medskip
\noindent Fie $K'\leq N_G(H)/H$, $K'=\langle\hat a\rangle$,
$|K'|=p$.

\[
\varphi:N_G(H)\longrightarrow\frac{N_G(H)}{H},
\qquad \varphi(a)=\hat a.
\]
\[
\varphi\bigl(\varphi^{-1}(K')\bigr)=K'.
\]
% Demonstratia se continua pe pagina urmatoare a manuscrisului.
.
% Fara preambul si fara intrerupere fortata de pagina.
% Continuarea demonstratiei lemei 2 de pe pagina 7.
% Nota de fidelitate: manuscrisul pare sa aiba "a in A_0",
% desi A_0 este o submultime de clase aH. Lectura este pastrata.

\subsection*{Curs 10 - Pagina 8}

\noindent $a\in A_0\ \Longleftrightarrow\ haH=aH$, $(\forall)\,h\in H$
\[
\Longleftrightarrow\ a^{-1}haH=H,\quad(\forall)\,h\in H
\]
\[
\Longleftrightarrow\ a^{-1}Ha=H
\Longleftrightarrow\ H=aHa^{-1}
\Longleftrightarrow\ a\in N_G(H).
\]

\noindent Deci
\[
|A_0|=\left|\frac{N_G(H)}{H}\right|=[N_G(H):H],
\]
\noindent iar $|A|=[G:H]$.

\bigskip
\noindent\underline{Corolar}. Aceleaşi ipoteze ca la lema 2 şi
$p\mid[G:H]$.

\noindent Atunci $H\subsetneq N_G(H)$.

\bigskip
\noindent\fbox{\textbf{Dem. T. lui Sylow 1}}\quad $p^i\mid |G|$.

\medskip
\noindent Inducţie după $i$.

\noindent $i=1$: O.K. din T. lui Cauchy.

\medskip
\noindent Presupunem $1\leq i<t$ şi că există $H\leq G$ a.î.
$|H|=p^i$.

\noindent Demonstrăm că există $K\leq G$, $|K|=p^{i+1}$,
şi $H\trianglelefteq K$.

\medskip
\noindent $p\mid[G:H]$, deoarece
\[
|G|=p^t m=[G:H]\,|H|=[G:H]p^i.
\]
\noindent $i<t\ \Longrightarrow\ p^{t-i}\mid[G:H]
\ \Longrightarrow\ p\mid[G:H]$.

\medskip
\noindent Din corolar, $H\subsetneq N_G(H)$.
Considerăm $N_G(H)/H$.

\medskip
\noindent Afirm că
\[
p\mid\left|\frac{N_G(H)}{H}\right|=[N_G(H):H],
\]
\noindent deoarece $p\mid[G:H]$ şi
\[
[N_G(H):H]\equiv[G:H]\pmod p.
\]

\noindent Din T. lui Cauchy, $(\exists)\,\hat a\in N_G(H)/H$
a.î. $\operatorname{ord}(\hat a)=p$.

\medskip
\noindent Fie $K'\leq N_G(H)/H$, $K'=\langle\hat a\rangle$,
$|K'|=p$.

\[
\varphi:N_G(H)\longrightarrow\frac{N_G(H)}{H},
\qquad \varphi(a)=\hat a.
\]
\[
\varphi\bigl(\varphi^{-1}(K')\bigr)=K'.
\]
% Demonstratia se continua pe pagina urmatoare a manuscrisului.
.
% Fara preambul si fara intrerupere fortata de pagina.
% Continuarea demonstratiei lemei 2 de pe pagina 7.
% Nota de fidelitate: manuscrisul pare sa aiba "a in A_0",
% desi A_0 este o submultime de clase aH. Lectura este pastrata.

\subsection*{Curs 10 - Pagina 8}

\noindent $a\in A_0\ \Longleftrightarrow\ haH=aH$, $(\forall)\,h\in H$
\[
\Longleftrightarrow\ a^{-1}haH=H,\quad(\forall)\,h\in H
\]
\[
\Longleftrightarrow\ a^{-1}Ha=H
\Longleftrightarrow\ H=aHa^{-1}
\Longleftrightarrow\ a\in N_G(H).
\]

\noindent Deci
\[
|A_0|=\left|\frac{N_G(H)}{H}\right|=[N_G(H):H],
\]
\noindent iar $|A|=[G:H]$.

\bigskip
\noindent\underline{Corolar}. Aceleaşi ipoteze ca la lema 2 şi
$p\mid[G:H]$.

\noindent Atunci $H\subsetneq N_G(H)$.

\bigskip
\noindent\fbox{\textbf{Dem. T. lui Sylow 1}}\quad $p^i\mid |G|$.

\medskip
\noindent Inducţie după $i$.

\noindent $i=1$: O.K. din T. lui Cauchy.

\medskip
\noindent Presupunem $1\leq i<t$ şi că există $H\leq G$ a.î.
$|H|=p^i$.

\noindent Demonstrăm că există $K\leq G$, $|K|=p^{i+1}$,
şi $H\trianglelefteq K$.

\medskip
\noindent $p\mid[G:H]$, deoarece
\[
|G|=p^t m=[G:H]\,|H|=[G:H]p^i.
\]
\noindent $i<t\ \Longrightarrow\ p^{t-i}\mid[G:H]
\ \Longrightarrow\ p\mid[G:H]$.

\medskip
\noindent Din corolar, $H\subsetneq N_G(H)$.
Considerăm $N_G(H)/H$.

\medskip
\noindent Afirm că
\[
p\mid\left|\frac{N_G(H)}{H}\right|=[N_G(H):H],
\]
\noindent deoarece $p\mid[G:H]$ şi
\[
[N_G(H):H]\equiv[G:H]\pmod p.
\]

\noindent Din T. lui Cauchy, $(\exists)\,\hat a\in N_G(H)/H$
a.î. $\operatorname{ord}(\hat a)=p$.

\medskip
\noindent Fie $K'\leq N_G(H)/H$, $K'=\langle\hat a\rangle$,
$|K'|=p$.

\[
\varphi:N_G(H)\longrightarrow\frac{N_G(H)}{H},
\qquad \varphi(a)=\hat a.
\]
\[
\varphi\bigl(\varphi^{-1}(K')\bigr)=K'.
\]
% Demonstratia se continua pe pagina urmatoare a manuscrisului.
.
% Fara preambul si fara intrerupere fortata de pagina.
% Continuarea demonstratiei lemei 2 de pe pagina 7.
% Nota de fidelitate: manuscrisul pare sa aiba "a in A_0",
% desi A_0 este o submultime de clase aH. Lectura este pastrata.

\subsection*{Curs 10 - Pagina 8}

\noindent $a\in A_0\ \Longleftrightarrow\ haH=aH$, $(\forall)\,h\in H$
\[
\Longleftrightarrow\ a^{-1}haH=H,\quad(\forall)\,h\in H
\]
\[
\Longleftrightarrow\ a^{-1}Ha=H
\Longleftrightarrow\ H=aHa^{-1}
\Longleftrightarrow\ a\in N_G(H).
\]

\noindent Deci
\[
|A_0|=\left|\frac{N_G(H)}{H}\right|=[N_G(H):H],
\]
\noindent iar $|A|=[G:H]$.

\bigskip
\noindent\underline{Corolar}. Aceleaşi ipoteze ca la lema 2 şi
$p\mid[G:H]$.

\noindent Atunci $H\subsetneq N_G(H)$.

\bigskip
\noindent\fbox{\textbf{Dem. T. lui Sylow 1}}\quad $p^i\mid |G|$.

\medskip
\noindent Inducţie după $i$.

\noindent $i=1$: O.K. din T. lui Cauchy.

\medskip
\noindent Presupunem $1\leq i<t$ şi că există $H\leq G$ a.î.
$|H|=p^i$.

\noindent Demonstrăm că există $K\leq G$, $|K|=p^{i+1}$,
şi $H\trianglelefteq K$.

\medskip
\noindent $p\mid[G:H]$, deoarece
\[
|G|=p^t m=[G:H]\,|H|=[G:H]p^i.
\]
\noindent $i<t\ \Longrightarrow\ p^{t-i}\mid[G:H]
\ \Longrightarrow\ p\mid[G:H]$.

\medskip
\noindent Din corolar, $H\subsetneq N_G(H)$.
Considerăm $N_G(H)/H$.

\medskip
\noindent Afirm că
\[
p\mid\left|\frac{N_G(H)}{H}\right|=[N_G(H):H],
\]
\noindent deoarece $p\mid[G:H]$ şi
\[
[N_G(H):H]\equiv[G:H]\pmod p.
\]

\noindent Din T. lui Cauchy, $(\exists)\,\hat a\in N_G(H)/H$
a.î. $\operatorname{ord}(\hat a)=p$.

\medskip
\noindent Fie $K'\leq N_G(H)/H$, $K'=\langle\hat a\rangle$,
$|K'|=p$.

\[
\varphi:N_G(H)\longrightarrow\frac{N_G(H)}{H},
\qquad \varphi(a)=\hat a.
\]
\[
\varphi\bigl(\varphi^{-1}(K')\bigr)=K'.
\]
% Demonstratia se continua pe pagina urmatoare a manuscrisului.
